%======================================================================
\subsection{BW-ROU-08 plan: flexible wires that keep their length (30 Oct 2026)}\label{sec:rou-flex-plan}
%======================================================================

\paragraph{The request.} You asked, through the session building the battery (\code{STD-02-01001-001}: two
12\,AWG leads from the pack to a T plug, four 22\,AWG balance leads to an XHP-4), for wires that behave like
real wires. Today a wire is a spline through its path points. Dragging the plug re-routes it only on release,
and the spline stretches (81\,mm became 132\,mm in the test). Your answers, already given:
\begin{enumerate}[nosep]
  \item the wire \textbf{keeps its length} (its cut length) and never bends tighter than its minimum radius;
        the ends stay clamped to the pin frames and leave them straight, as now;
  \item it moves \textbf{live during an Assembly drag}, solved again at each drag step from the last shape;
  \item pulled taut, \textbf{the drag stops} at the wire's length, like the joint-limit drag stop of
        \S\ref{sec:joints} (no warn and continue);
  \item \textbf{gravity sag} is in;
  \item \textbf{collisions with parts} are in, against simplified shapes for speed;
  \item this is the \textbf{default for every wire}, not a mode.
\end{enumerate}
The method borrows the \emph{ideas} of Blender's hair and cloth tools, not their code (Blender is GPL): position-based
dynamics for the rope and discrete elastic rods for the stiffness. Everything is numpy; nothing new is bundled.

%----------------------------------------------------------------------
\subsubsection*{Theory}

\paragraph{The rope.} A wire of cut length $L$ is $N$ points $\mathbf x_0\dots\mathbf x_{N-1}$ with a rest
segment length $\ell_0=L/(N-1)$. The first and last $k$ points sit on the leads: they are fixed on the pin
frames' $+Z$ axes for the straight lengths of ROU-05, so the wire leaves each terminal straight. The rest move.
The shape is the \emph{resting} configuration under three kinds of constraint and gravity. It is not a dynamic
simulation: we want where the wire settles, not how it swings.

\paragraph{XPBD (Macklin, M\"uller, Chentanez 2016).} Extended position-based dynamics moves points directly to
satisfy constraints $C_j(\mathbf x)=0$ (or $\ge0$), each with a compliance $\alpha_j$ (the inverse of a
stiffness). One step of length $\Delta t$:
\[
  \tilde{\mathbf x}_i=\mathbf x_i+\Delta t\,\mathbf v_i+\Delta t^2\,\mathbf g,\qquad
  \Delta\lambda_j=\frac{-C_j-\tilde\alpha_j\lambda_j}{\sum_i w_i\lVert\nabla_iC_j\rVert^2+\tilde\alpha_j},\qquad
  \Delta\mathbf x_i=w_i\,\nabla_iC_j\,\Delta\lambda_j,
\]
with $\tilde\alpha_j=\alpha_j/\Delta t^2$, inverse masses $w_i$ (zero for the clamped points), and the
Lagrange multipliers $\lambda_j$ accumulated over the iterations. With heavy velocity damping and the previous
shape as the start, a few steps reach rest. That is what lets a drag step cost milliseconds, because each step
starts from where the wire already is.

\paragraph{Constraints.}
\begin{itemize}
  \item \textbf{Length} (inextensible, $\alpha=0$): $C=\lVert\mathbf x_{i+1}-\mathbf x_i\rVert-\ell_0$ per segment.
        The wire neither stretches nor shrinks; its length is $L$ by construction, not by checking.
  \item \textbf{Bending} (elastic): the turning angle $\theta_i$ at point $i$ between consecutive segments.
        The discrete elastic rod (Bergou et al.\ 2008) gives the bending energy
        $E_b=\sum_i \frac{EI}{\ell_0}\bigl(2\tan\frac{\theta_i}{2}\bigr)^2$, so the compliance is
        $\alpha_b=\ell_0/EI$. Twist is left out: a round wire's twist barely changes its centreline.
  \item \textbf{Minimum bend radius} (a limit, not a spring): the local radius $R_i=\ell_0/(2\sin\frac{\theta_i}{2})$
        must stay at or above $R_{\min}$ from the wire's bend rule, i.e.\ $\theta_i\le\theta_{\max}=2\arcsin
        \frac{\ell_0}{2R_{\min}}$. It is active only when violated. For 12\,AWG at $4\times$OD ($R_{\min}=14.4$\,mm)
        and $\ell_0=6$\,mm, $\theta_{\max}=24^\circ$.
  \item \textbf{Collision} (inequality): for each point and each nearby part proxy with signed distance $d(\mathbf x)$,
        $C=d(\mathbf x_i)-r\ge0$ with $r$ the wire radius plus a clearance. A violated point is pushed out along
        $\nabla d$.
\end{itemize}

\paragraph{Stiffness from the gauge.} A stranded conductor bends with its strands sliding, so its stiffness is
close to the sum of the strands', $EI_c\approx n_s E_{\text{Cu}}\pi d_s^4/64$, a lower bound (friction adds
some). The insulation is a tube, $EI_i=E_i\pi(D^4-d^4)/64$. For 12\,AWG (65 strands of 0.254\,mm, PVC OD
3.6\,mm): $EI_c\approx1554$, $EI_i\approx137$, $EI\approx1.7\times10^3$\,N\,mm$^2$. A table per insulation type
(PVC, silicone, PTFE, XLPE, matching the ROU bend-rule presets) holds $E_i$ and density, and a per-wire factor
calibrates it against a hang test (question W9).

\paragraph{Gravity.} The weight per length $q=(\rho_{\text{Cu}}A_c+\rho_iA_i)g$ is
$3.7\times10^{-4}$\,N/mm for 12\,AWG. Gravity points along $-Z$ of the top assembly unless the assembly says
otherwise (question W7).

\paragraph{Taut and the drag stop.} An inextensible rope whose end points move apart beyond its length has no
solution. The solver detects this as a residual: the length constraints cannot all be met, and the largest
segment strain stays above $10^{-3}$ after the iterations. The quick test first is the lower bound
$\lVert\mathbf a-\mathbf b\rVert+\ell_a+\ell_b\le L$ between the lead tips $\mathbf a,\mathbf b$ (with obstacles
the shortest path is longer, so the residual decides). The drag step is then cut back by bisection, as the joint
limits already are (\S\ref{sec:joints}): the part stops where the wire just reaches its length.

\begin{figure}[htbp]
\centering
\begin{tikzpicture}[x=1mm,y=1mm,font=\sffamily\scriptsize]
  \fill[black!8] (0,10) rectangle (12,34); \node at (6,36) {pack};
  \fill[black!8] (108,16) rectangle (120,40); \node at (114,42) {plug};
  \fill[bworange!25] (48,0) rectangle (70,12); \node at (59,6) {part (proxy box)};
  \draw[black!40, dashed] (12,22) -- (20,22); \draw[black!40, dashed] (100,28) -- (108,28);
  \draw[bwblue, line width=1.1pt] plot[smooth, tension=0.6] coordinates {(12,22) (20,22) (32,18) (46,15) (59,14.5) (72,16) (86,21) (100,28) (108,28)};
  \foreach \p in {(12,22),(16,22),(20,22),(108,28),(104,28),(100,28)} \fill[black] \p circle (0.8);
  \foreach \p in {(26,20),(32,18),(39,16.3),(46,15),(52.5,14.6),(59,14.5),(65.5,15),(72,16),(79,18.3),(86,21),(93,24.5)} \draw[bwblue, fill=white] \p circle (0.8);
  \draw[->, black!50] (5,4) -- (5,-2) node[right] {$\mathbf g$};
  \node[anchor=west, text=black!55] at (2,46) {$\bullet$ clamped on the lead (straight length of the pin)};
  \node[anchor=west, text=black!55] at (2,50) {$\circ$ free points: length, bending, bend limit, collision, gravity};
  \node[anchor=west, text=black!55] at (60,50) {the part proxy pushes points out by $r$};
\end{tikzpicture}
\caption{A wire as a rope of points. The leads are clamped along the pin frames; the free points settle under
gravity while every segment keeps $\ell_0$, every bend stays above $R_{\min}$, and no point enters a part proxy.}
\label{fig:rou-flex}
\end{figure}

%----------------------------------------------------------------------
\subsubsection*{How it fits Keystone}

\begin{itemize}
  \item \textbf{The wire keeps a state.} A new hidden property stores the rope's points. A recompute that moves a
        pin (an edited joint, a moved part) solves from that state, so a wire never jumps to a new shape.
  \item \textbf{During a drag.} Assembly drags go through \code{preDrag}/\code{doDragStep}/\code{postDrag}
        (commit 1a15d0d1). A small \code{[BW]} hook in \code{AssemblyObject::doDragStep} calls the registered
        Python \emph{drag checks} after the solver proposes placements and before they are applied. The rope
        module is one of them: it solves the wires attached to the moved parts, returns ``ok'' or ``taut'', and
        a taut answer sends the step into the existing bisection. While dragging, a wire is drawn as a light
        tube (a Coin line set with the wire's width and colour), because sweeping an OCC solid 30 times a second is
        too slow; the real solid is built once on release.
  \item \textbf{Collision proxies.} Each part near a wire (bounding boxes within the wire's reach) gets a proxy:
        its oriented bounding box by default, or a convex hull or decimated mesh for parts where the box is too coarse
        (question W6). Signed distances to boxes are exact and cheap; to a mesh, through a small BVH in numpy.
  \item \textbf{Resolution.} About one point per 1.5\,OD, clamped to 20--200 points (question W4); 40 for the
        battery's 81\,mm lead.
  \item \textbf{Flexible sub-assemblies.} An \code{AssemblyLink} with \code{Rigid=False} copies only links and
        assembly links into the parent (\code{synchronizeComponents}). Wires and route paths stay behind in the
        sub-assembly's document and do not follow parts the parent moves. The fix: the copy step also copies
        routed items, re-pointing their attachments to the copied links (C++, \code{[BW]}).
\end{itemize}

\paragraph{Verification.}
\begin{center}
\small
\begin{tabularx}{\linewidth}{@{}lX@{}}
\toprule
Catenary & $EI\to0$, ends 200\,mm apart at the same height, $L=250$\,mm: the shape is a catenary with
  $a=84.55$\,mm and sag $a(\cosh\frac{d}{2a}-1)=66.36$\,mm; the test accepts 1\,\%.\\
Cantilever stub & 12\,AWG, 100\,mm clamped horizontally at one end, free at the other: tip sag
  $qL^4/8EI=2.76$\,mm (small deflection); accepts 5\,\%.\\
Length & any drag sequence: $\sum\lVert\mathbf x_{i+1}-\mathbf x_i\rVert=L$ within $10^{-3}$.\\
Bend limit & a pin turned back on itself: no turning angle above $\theta_{\max}$.\\
Drag stop & the T plug dragged away from the pack stops where the 12\,AWG leads are taut, within one bisection
  step; dragging back frees it.\\
Collision & a wire over a box sags onto it and no point is inside the box minus $10^{-2}$\,mm.\\
Battery & \code{STD-02-01001-001} (on a vault copy): drag P1 through an arc; the leads keep 81\,mm, update live
  ($\ge20$ frames/s), stop when taut.\\
\bottomrule
\end{tabularx}
\end{center}

\paragraph{New epic.} \eref{ROU-08} Flexible wires: rope solver as the default wire behaviour, live in drags, drag
stop when taut, gravity, collision proxies, flexible sub-assemblies carrying routed items.

\paragraph{Questions for you.}
\begin{enumerate}[label=W\arabic*.]
  \item \textbf{Hand-placed path points.} Under the rope, what is a point you picked? \emph{Recommendation:} a
        point picked on a part (a clip, an LCS) becomes a \emph{clip} the wire passes through, free to turn; a point
        in free space is only a \emph{guide} for the first shape and then forgotten.
  \item \textbf{Migrating existing wires.} Each wire's cut length becomes its current length (with its slack), so
        nothing changes shape until something moves? \emph{Recommendation:} yes, and the slack field then edits the
        cut length directly.
  \item \textbf{Bundles, hoses, pipes.} \emph{Recommendation:} hoses follow (same rope, hose stiffness); a bundle is
        one rope with the combined OD and stiffness, branches at breakouts later; pipes and tubes stay rigid (they are
        bent, not flexible).
  \item \textbf{Resolution.} One point per 1.5\,OD, 20--200 points? \emph{Recommendation:} yes.
  \item \textbf{Live display.} A light tube while dragging, the real solid on release? \emph{Recommendation:} yes.
  \item \textbf{Collision proxies.} Oriented box per part by default, convex hull or mesh chosen per part; wire to
        wire contact later? \emph{Recommendation:} yes.
  \item \textbf{Gravity.} $-Z$ of the top assembly, settable per assembly (a drone shown on its side)?
        \emph{Recommendation:} yes.
  \item \textbf{Outside drags.} A recompute that moves a pin (joint edit, typed placement) solves the rope from its
        last shape, and a taut wire there turns red with a badge rather than refusing the edit?
        \emph{Recommendation:} yes. Only drags can be stopped.
  \item \textbf{Stiffness data.} Table per insulation type plus a per-wire calibration factor from a hang test (clamp
        100\,mm, measure the tip sag)? \emph{Recommendation:} yes.
  \item \textbf{Bend limit while dragging.} If a drag would force a bend tighter than the minimum (the pin turned hard),
        stop the drag there too, or let the wire turn red? \emph{Recommendation:} stop the drag, like taut.
  \item \textbf{Flexible sub-assemblies.} Copy routed items into the parent with the links (above)?
        \emph{Recommendation:} yes.
\end{enumerate}

\paragraph{Your answers (30 Oct).} All recommendations W1--W11 accepted, with three additions:
\begin{itemize}[nosep]
  \item \textbf{Any length.} A wire or hose can be made at any cut length you type, longer than its route (it
        sags, loops or lies on the parts below) or as short as its leads allow (then it is taut, and red outside a
        drag). The cut length is the wire's own property, not derived from the path; the path's points only give
        the first layout (W1).
  \item \textbf{Existing wires do not extend}: they keep their present length as their cut length (W2), and no
        wire ever stretches.
  \item \textbf{Gravity is always on}, and its direction is chosen per assembly along an axis ($\pm X$, $\pm Y$,
        $\pm Z$ of the assembly, default $-Z$), W7.
\end{itemize}
Also accepted: Insert Fastener files new bolts under a \emph{Hardware} category (\S\ref{sec:categories-built}).
